\documentclass[11pt,a4paper]{article}
% Generated from papers/research-release.json.
\newcommand{\PublicationVersion}{DR0149-r3}
\newcommand{\PublicationDate}{9 November 2026}
\newcommand{\PublicationPart}{Part I: Research Overview and Certified Recognition}

\usepackage[T1]{fontenc}
\usepackage[utf8]{inputenc}
\usepackage{lmodern,amsmath,amssymb,amsthm,booktabs,longtable,array,geometry,microtype}
\geometry{margin=27mm}
\usepackage{xurl}
\usepackage[colorlinks=true,linkcolor=blue,citecolor=blue,urlcolor=blue]{hyperref}
\hypersetup{pdftitle={Nullivance Logic: Foundations, Quantified Reasoning, and Certified Recognition},pdfauthor={Trinh Tung Lam},pdfsubject={Part I: Research Overview and Certified Recognition; DR0149-r3}}
\setlength{\emergencystretch}{3em}
\newtheorem{theorem}{Theorem}[section]
\newtheorem{proposition}[theorem]{Proposition}
\newtheorem{definition}[theorem]{Definition}
\newtheorem{example}[theorem]{Example}
\newtheorem{remark}[theorem]{Remark}
\newcommand{\leanref}[1]{\par\smallskip\noindent{\footnotesize Lean anchor: \nolinkurl{#1}.}\par\smallskip}
\newcommand{\dr}[1]{DR-#1}
\newcommand{\Four}{\mathbf{4}}
\newcommand{\Qv}{\mathsf{Q}}
\newcommand{\Fits}{\mathsf{Fits}}
\newcommand{\Known}{\mathsf{Known}}
\newcommand{\none}{\mathsf{none}}
\newcommand{\some}{\mathsf{some}}
\title{Nullivance Logic:\\Foundations, Quantified Reasoning,\\and Certified Recognition}
\author{Trinh Tung Lam\\\small\href{mailto:lamtung0481@gmail.com}{lamtung0481@gmail.com}}
\date{\PublicationDate\\\small \PublicationPart\ --- \PublicationVersion}
\begin{document}
\maketitle
\begin{abstract}
When distinct latent states produce the same observations, correct reasoning
about those observations need not identify the latent state. We study this
separation in Nullivance, a two-channel logic with an optional intensity--structure
model. The logical foundation admits exact threshold projection and a complete
signed tableau; quantified results cover a finite-domain calculus and a specified
pure-equality fragment on finite and infinite carriers. For recognition, we prove
that the initial observation cannot determine the defined quasivance predicate.
An explicitly stronger scalar probe interface permits exact classification of
finite rational histories into infeasible, affirmed, refuted, or undetermined
cases, with witnesses and source-backed certificates. We then distinguish
compatibility of an infinite history from certification after finite observation.
At a positive intensity threshold equal to the latent maximum intensity, finite
resolution occurs exactly when an observed lower endpoint contacts the threshold.
An unknown period does not suffice for a universally conclusive finite negative
procedure; a supplied certified period bound $B$ yields a scanner with at most
$B$ contact tests and exact selected-operation counts. Lean 4 checks the stated
formal results and implementation refinements. The contribution is a scoped,
traceable development connecting semantics, observation limits, and certified
algorithms, not a claim to have invented four-valued reasoning or solved general
cognition. No empirical performance or cognitive experiment is reported.
\end{abstract}
\noindent\textbf{Keywords:} non-classical logic; formal verification; equality;
observability; recognition; certificates; Lean 4.
\par\medskip
\noindent\textbf{Reading this collection.} This overview connects the project
and develops its recognition results. Part II gives the propositional calculus;
Part III gives the finite-domain quantified calculus. All three documents carry
collection revision \PublicationVersion. Archived earlier editions remain
available; they are not interchangeable with this edition. The source bundle
contains the Lean development, reproducible builds, and a version history.
\clearpage
\tableofcontents
\newpage

\section{Introduction}
\label{sec:scope}
\subsection{The problem: reasoning from reports is not recovering a state}
Suppose two instruments disagree about whether a component is active. One
may record support for activity without discarding support against it. A logic
can specify exactly which conclusions follow from those reports. That task is
different from deciding whether the component is actually active: the latter
also depends on how the instruments are related to the component. A correct
proof system cannot recover information that its observation interface omits.
This distinction motivates the present work. The instrument example is an
illustration, not an experiment or a validated physical model.

The original Nullivance intuition asks how an unmanifested structure should be
distinguished from an absence of structure. We turn this question into a
mathematical one: given a latent state space, a target property, and an explicit
observation map, which differences between states survive observation? A zero
output need not identify its cause. Conversely, declaring a latent coordinate
in a model does not demonstrate the existence of a corresponding physical entity.
The term \emph{quasivance} below names one precise predicate, not a scientific
claim about entities acting without existing.

\subsection{Research questions and results}
The development addresses three connected questions.
\begin{enumerate}
\item \emph{Logical consequence.} Can independent positive and negative
support be given exact semantics and a complete proof system? The propositional
answer uses threshold projection onto four values and signed tableaux. The
quantified answer is restricted: a complete finite-domain calculus and exact
algorithms for pure equality, not unrestricted infinite-domain completeness.
\item \emph{Information and certificates.} Which latent properties can be
recovered from specified observations? The initial generative observation loses
the quasivance distinction. Under stronger scalar probes, compatible states
admit an exact box representation. Finite histories then support complete
classification and explicit witnesses, including a distinction between
inconsistent evidence and genuine uncertainty.
\item \emph{Finite stopping.} When does sufficient information become available
after finitely many observations? Infinite compatibility, identifiability from
an entire signature, and finite termination are different. Our boundary-contact
criterion supports a matched negative/positive pair: unknown periodicity is
insufficient for universal finite negative decisions, while a certified upper
bound on the contact period gives a correct bounded scanner.
\end{enumerate}
The implementation layer connects these specifications to memoized equality
evaluation, checked memory import, resumable replay, and operation counters.
The central contribution claimed here is this formally checked chain for the
specified models. Standard ingredients and elementary general principles are
identified as such; the number of development records is not a measure of novelty.

\subsection{Relation to established work}
The truth/falsity separation belongs to the Belnap--Dunn and bilattice tradition
\cite{belnap1977useful,dunn1976intuitive,fitting1991bilattices}. The connective
called harmonization is the coordinatewise minimum, i.e. bilattice consensus,
not a newly discovered algebraic operation. Part II gives the exact fragment
embeddings and literature comparison. Our contribution here is not a claim that
classical mathematics cannot describe uncertain or contradictory reports.

Epistemic accounts of distributed systems define knowledge through alternatives
compatible with local information \cite{halpern1990knowledge}. Our observation
fiber criterion is a single-observer factorization principle in that tradition.
It is not a new modal logic of multi-agent common knowledge. The present target
is the specific alpha/Theta predicate and its executable finite evidence contracts.

Representing states consistent with bounded-error data is a standard
set-membership estimation problem. Jaulin and Walter's interval approach
handles nonlinear set inversion by guaranteed enclosures
\cite{jaulin1993inversion}. Here the enriched probe constraints are rectangular
and coordinatewise, so intersection is exact; we do not claim a solution of
general nonlinear identification or superiority to interval-estimation methods.

The sufficient syntactic analyses of regularity are related in method to
sound abstraction \cite{cousot1977abstract}; failure to certify a property is
not its refutation. Ordered reduced Boolean diagrams are due to established
work \cite{bryant1986graph}. Our development proves contracts for its particular
representations and combinations. This comparison positions the results but
is not an exhaustive priority survey of all recognition or verification research.

\subsection{Method and organization}
We state a model and its input contract, prove a result or construct a
counterexample within that model, and link the statement to a Lean declaration.
Implementation refinements have separate correctness and cost claims.
The method is mathematical proof and formal verification, not experimental
confirmation of the model. Worked numerical cases illustrate definitions;
they do not replace universally quantified proofs.

Three different objects must be kept separate. The first is a logic of
independent support for and against propositions. The second is a generative
model for the values supplied to that logic. The third is an observation model
in which an observer attempts to recover specified properties of a latent
state. A theorem about one object is not automatically a theorem about the
others. In particular, evidence in both channels is permitted by the logic;
factive knowledge of a proposition and its negation is not.

The exposition proceeds from logic and quantification to the latent model,
finite observations, infinite histories, and stopping. Part II and Part III
supply the full tableau rule sets and their detailed proof arguments, so the
three-part collection should be used for those dependencies. This overview
is not a replacement for those calculi. Appendix~\ref{app:catalog} is a
provenance index, separate from the mathematical argument. Full Lean types
retain all semantic and typeclass hypotheses; resolving a theorem name alone
does not verify a natural-language paraphrase.

\begin{center}\small
\begin{tabular}{p{.26\linewidth}p{.64\linewidth}}
\toprule
Layer & Established scope; important exclusion\\\midrule
Propositional core & Signed tableau soundness/completeness and exact threshold
projection; no interpretation as a physical law.\\
Quantification & Finite-domain core tableau; finite/infinite pure-equality
algorithms; not arbitrary-predicate infinite-domain completeness.\\
Recognition & Exact properties of specified observations of a fixed scalar
state; not unrestricted perception, changing-state cognition, or learning.\\
Algorithms & Correctness and explicitly charged operations of named programs;
not a universal wall-clock, memory, or bit-complexity bound.\\
Publication evidence & Compiling proof artifact, dependency audits, source
manifest, and reproducibility instructions; not external peer review.\\
\bottomrule
\end{tabular}
\end{center}

Definitions stipulate a model; proofs establish consequences of that model.
Any application must independently justify its observation map, fixed-state
assumption, error bounds, and supplied stopping certificates. These are not
provided by the logical connectives or by the proof assistant.

\section{Two-channel semantics and the propositional core}
An atomic proposition receives a pair $(t,f)\in[0,1]^2$. These coordinates
are independent support values, not complementary probabilities. For values
$x=(t,f)$ and $y=(u,v)$, the primitive operations are
\[
 \neg x=(f,t),\quad x\wedge y=(\min(t,u),\max(f,v)),\quad
 x\vee y=(\max(t,u),\min(f,v)),
\]
and harmonization (consensus) is $x\oplus y=(\min(t,u),\min(f,v))$.
A threshold $0<\tau\le1$ records whether each coordinate is at least $\tau$.
The four resulting Boolean pairs are $\mathrm{N}=(0,0)$, $\mathrm{T}=(1,0)$,
$\mathrm{F}=(0,1)$, and $\mathrm{B}=(1,1)$. Here zero and one in a Boolean
pair denote bits. The signs $T^+,T^-,F^+,F^-$ assert, respectively, positive
truth support, its absence, positive falsity support, and its absence at the
specified threshold. Absence of support is not identified with support for
the opposite proposition.

\begin{example}
At $\tau=1/2$, $(0.8,0.7)$ projects to B, whereas $(0.2,0.1)$ projects to N.
The first can represent conflicting reports; the second insufficient reports.
Neither description says that an observer knows both a fact and its negation.
The same proposition can have different reports without changing the
mathematical meanings of the two channels.
\end{example}

\begin{theorem}[Exact propositional projection]
Thresholding the continuous evaluation of a formula gives exactly its
four-valued evaluation under thresholded atomic values.
\end{theorem}
\begin{proof}
Induct on the formula. Swapping coordinates commutes with thresholding.
For two reals, a minimum is at least $\tau$ iff both operands are; a maximum
is at least $\tau$ iff at least one is. These identities establish all three
binary cases coordinate by coordinate. This argument uses finite min/max,
not an unattained supremum of an infinite family.
\end{proof}
\leanref{Nullivance.Continuous.exact_projection}

The analytic signed tableau closes precisely the unsatisfiable finite signed
branches. The corresponding consequence theorem transfers to the continuous
semantics by exact projection. The artifact also supplies executable reference
checking, countermodel extraction on failure, and compactness/strong-completeness
results for infinite premise sets. The last phrase concerns \emph{sets of
propositional premises}; it does not mean arbitrary quantified object domains.
\leanref{Nullivance.Metatheory.Closes.unsat}
\leanref{Nullivance.Metatheory.closes_of_unsat}
\leanref{Nullivance.Metatheory.derives_iff_consequenceC}

Paraconsistency means that contradictory support need not entail an arbitrary
unrelated proposition. It does not make every inference permissible. The
two companion papers contain the detailed core calculus and finite-domain
tableau exposition. Part II distinguishes the paper-level coNP classification
from its machine-checked certificate and reduction components; its FDE and
bilattice connections have explicit formal anchors. These statuses are not
interchangeable. The mathematical core deliberately builds on established
four-valued and bilattice ideas \cite{belnap1977useful,dunn1976intuitive,fitting1991bilattices}.
The availability of a Lean formalization alone does not establish historical
novelty of those ideas.

\section{Quantifiers: finite results, infinite obstructions, and equality}
\subsection{Finite-domain proof theory}
The finite-domain work treats a function-free quantified language with crisp
equality and a nonempty finite domain. The macro-free extensional core tableau
has a soundness and completeness theorem. Equality substitution and replay
require their actual admissibility conditions; earlier proposed replay routes
were corrected rather than left as assumed lemmas. Part III distinguishes those failed routes from the
proved final calculus. A refuted auxiliary conjecture does not refute the
subsequently repaired theorem.
\leanref{Nullivance.FiniteFO.qDerivesExtCore_iff_qconsequence4Sig}

\subsection{Why finite projection is insufficient for all infinite domains}
Consider truth values $a_n<\tau$ increasing to $\tau$. The existential truth
coordinate may be $\sup_n a_n=\tau$, so its continuous threshold bit is true,
although every individual threshold bit is false. Moving projection through
this supremum is invalid. Positive results therefore require stated witness
attainment, threshold regularity, or suitable compactness and semicontinuity
hypotheses. Equality depending on a bound variable also matters: an arbitrary
crisp equality predicate need not have the topological regularity that an
informal continuity argument assigns to it.

The abstract analyses track constant, crisp, and upper-semicontinuous
coordinates, then coordinate relations and exact complementarity. These are
sound sufficient recognizers, not complete recognizers for every semantic
property. Boolean normalization by ROBDD removes a particular syntactic
limitation: equivalent Boolean combinations need not have the same syntax
tree. But a propositional ROBDD treats equality atoms as independent unless
the equality theory is added. For example, reflexivity of $x=x$ and transitive
consequences of $x=y,y=z$ are theory facts, not arbitrary Boolean tautologies.
Reduced ordered decision diagrams are an established Boolean representation
\cite{bryant1986graph}; the present artifact verifies their use in its own
pipeline, not their invention.

\subsection{Exact equality procedures}
The equality branch adds equivalence classes, disequality constraints,
binder-aware assignment updates, and fresh-element capacity. Universal and
existential choices are represented by existing classes and, when possible,
a fresh class. Forgetting which variables are currently bound can merge
semantically different calls; the cache keys therefore retain the necessary
scope/capacity information. A cutoff argument replaces infinitely many
indistinguishable choices by finitely many equality types.

The resulting total decision procedures are exact for the stated
\emph{pure-equality fragment}. For closed formulas the artifact proves
correctness on infinite carriers and on nonempty finite carriers with the
appropriate capacity. It does not decide arbitrary predicates on an infinite
domain. The later open-formula interfaces evaluate a \emph{supplied assignment},
not automatically every assignment.
\leanref{Nullivance.InfiniteFO.decideQuantifiedEqualityHashed_infinite_correct}
\leanref{Nullivance.InfiniteFO.decideQuantifiedEqualityHashed_finite_correct}

For an open formula of quantifier rank $r$, let $c$ be the number of distinct
values assigned to its relevant free variables. Compact encodings use their
equality classes and an explicit fresh pool. The sufficient finite capacity
condition is $r+c\le |D|$; when it fails, the checked finite interface can
use a supplied complete duplicate-free enumeration of the actual domain.
The validation of the enumeration and fragment is part of the interface.
Rejection of an unsupported input is not a false truth value.
The conservative compact representative construction uses $r+c+1$ elements;
no minimality of that count is asserted here. Enumerating validity across all
free assignments remains separate from evaluating one assignment.
\leanref{Nullivance.InfiniteFO.decideFiniteOpenEquality_correct}
\leanref{Nullivance.InfiniteFO.checkedFiniteOpenEquality_sound}
\leanref{Nullivance.InfiniteFO.decideFreshCompactOpenEquality_finite_capacity_correct}

\subsection{Memoization and what the cost theorems measure}
The optimized evaluator is proved equivalent to its reference recursion;
the table is threaded through recursive orbit expansion. Its invariant
identifies distinct stored states with misses from the specified empty-table
execution. Syntax-sensitive Bell-number and factorial bounds then control
the number of unique keys. The exact bound parameters and closure conditions
are retained in the following source anchors, not suppressed into an
unqualified ``polynomial-time'' claim.
Writing $r$ for quantifier rank and $|\varphi|$ for syntax-node count, the
empty-root specialization proves
\[
 U=\mathsf{misses}\le |\varphi|\,B_r\le |\varphi|\,r!,
\]
where $B_r$ is the artifact's equality-partition (Bell) count at that root.
The stronger depth-sensitive envelope is available in the source; the simple
display is an upper bound, not the exact count of every run.
\leanref{Nullivance.InfiniteFO.runQuantifiedEqualityOrbitHashed_correct}
\leanref{Nullivance.InfiniteFO.runQuantifiedEqualityOrbitHashed_unique_states_eq_misses}
\leanref{Nullivance.InfiniteFO.runQuantifiedEqualityOrbitHashed_unique_states_le_syntax_mul_bell}
\leanref{Nullivance.InfiniteFO.runQuantifiedEqualityOrbitHashed_unique_states_le_syntax_mul_factorial}

For the pinned Lean implementation, a bucket-doubling potential bounds
rehash work linearly in the number $U$ of inserted unique states. In the
instrumented model, rehashed entries are bounded by $3U$ and the counted
hash operations by $A+4U$, where $A$ is the counted request quantity in the
formal cost specification. Collision-chain comparison costs retain their
dependence on chain length; no random hash distribution is postulated.
The displayed amortized operation bounds hold for the instrumented root run;
the syntax/Bell envelope specializes the parameter to quantifier rank.
Semantic correctness additionally requires the stated fragment and closure
conditions. These costs exclude the bit cost of a large key,
rational arithmetic, allocation, and the entire surrounding ROBDD algorithm.
\leanref{Nullivance.InfiniteFO.runQuantifiedEqualityOrbitHashed_amortized_state_cost_bounds}

\section{The generative layer: intensity is not observation}
A generative frame has a positive structure dimension $d$, a stability map
$\Phi:[0,1]^d\to[0,1]$, flip symmetry
$\Phi(1-\Theta)=\Phi(\Theta)$, and neutral value
$\Phi(1/2,\ldots,1/2)=1$. A channel contains an intensity $\alpha\in[0,1]$
and a structure $\Theta\in[0,1]^d$. Its effective output is
$\alpha\Phi(\Theta)$. A state has two channels and initializes the core value as
\[
 \mathsf{init}(s)=
 (\alpha_+\Phi(\Theta_+),\alpha_-\Phi(\Theta_-)).
\]
One concrete admissible family uses the geometric mean of componentwise
$1-2|\theta-1/2|$. This is a model definition, not a measured biological law.
\leanref{Nullivance.Generative.GenState.init_mem}

\begin{definition}[The specified latent predicate]
A channel is quasivant when $\alpha=0$ and its structure is not neutral.
A state is quasivant, written $\Qv(s)$, when both channels are quasivant.
\end{definition}
This exact conjunction is important. Small positive intensity is not zero;
a neutral structure is not nonneutral; one quasivant channel is not two.
All quasivant states initialize to $(0,0)$, but the converse fails.

\begin{proposition}[Silent ambiguity in every admitted frame]
There are a quasivant state and a nonquasivant state with identical initial
observations, for every admitted generative frame.
\end{proposition}
\begin{proof}
Set both intensities to zero. In one state set both structures constantly to
zero, and in the other constantly to $1/2$. Positive dimension ensures these
structures differ. The first state satisfies the defined predicate and the
second does not. Multiplication by zero makes both observations $(0,0)$,
independently of the admitted stability map.
\end{proof}
\leanref{Nullivance.Recognition.every_frame_silent_ambiguity}
This refutes universal recovery of this latent predicate from init alone.
It neither disproves the existence of structured latent states nor requires
changing the logical core.

\section{Recognition as recoverability from observations}
Let $o:S\to O$ be an observation map and $P:S\to\mathsf{Prop}$ a property.
Exact recognizability means that a decoder $d:O\to\mathsf{Prop}$ exists with
$d(o(s))\leftrightarrow P(s)$ for every state.
\begin{theorem}[Fiber criterion]
Such a decoder exists iff $P$ is constant on each observation fiber:
\[
 o(s)=o(t)\ \Longrightarrow\ (P(s)\leftrightarrow P(t)).
\]
\end{theorem}
\begin{proof}
A decoder gives the forward implication immediately. Conversely define
$d(y)$ by the existence of a state $s$ with $o(s)=y$ and $P(s)$. At any
actual observation $o(t)$, fiber constancy makes this equivalent to $P(t)$.
This constructs a mathematical predicate, not necessarily a computable test.
\end{proof}
\leanref{Nullivance.Recognition.recognizable_iff_fiber_constant}

Postprocessing cannot restore distinctions already removed by the observation.
Extra observations can preserve or improve recoverability. A family of probes
supports recognition iff it separates every pair disagreeing on the target
property; observing an entire possibly infinite signature is not the same
as finding a finite stopping rule.
\leanref{Nullivance.Recognition.recognizable_of_postprocess}
\leanref{Nullivance.Recognition.recognizable_tests_iff_separates}

If deterministic dynamics preserve observational equality at each next step,
induction preserves it along the whole trace. Merely appending a time axis
does not then recover the lost distinction. This is a conditional result about
observation-stable dynamics, not an impossibility for all interventions.
A coordinate-swap example shows how an intervention can reveal a hidden bit
when that stability condition is absent.
\leanref{Nullivance.Recognition.stable_evolution_indistinguishable}

Factive knowledge is separately specified by
$\Known_o(P,s)\leftrightarrow\forall t,\ o(t)=o(s)\Rightarrow P(t)$.
The actual state is a compatible alternative, so knowledge implies truth.
Four-valued support is not automatically knowledge. Even T alone is
insufficient without a sound connection to the target, and B cannot be
promoted factively to both a proposition and its negation. A positive bridge
requires adequacy of reports and constancy on the observation fiber.
\leanref{Nullivance.Recognition.truthSupport_known_of_adequate}
\leanref{Nullivance.Recognition.bothSupport_no_unconditional_knowledge}
These hypotheses are application obligations, not facts silently supplied by
the Nullivance connectives.

\section{Scalar probes, noise, and exact finite histories}
\label{sec:histories}
The constructive recognition branch specializes the generative frame to four
scalar coordinates $(\alpha_+,\alpha_-,\theta_+,\theta_-)\in[0,1]^4$ and
an enriched probe interface. Probe readings expose information absent from
init alone; this additional observation model is not derived by decoding init.
Concretely, the intensity intervention replaces structure by neutral and
retains intensity, so its response is exactly $\alpha$. In the canonical
one-dimensional frame, the structure intervention sets intensity to one and
halves the original structure coordinate; its response is exactly $\theta$.
These experiments operate on separate copies of the same original state.
Their combined exact-real signature is injective, hence information-theoretically
sufficient for any property of that state. The signature is not an executable
decision procedure for every predicate over real numbers, and the proposed
interventions have not been implemented as physical experiments.
\leanref{Nullivance.Recognition.activeProbeSignature_eq}
\leanref{Nullivance.Recognition.activeProbeSignature_injective}
The scalar target is
\[
 \Qv(s)\ \leftrightarrow\ \alpha_+=0\land\alpha_-=0
 \land\theta_+\ne1/2\land\theta_-\ne1/2.
\]
A rational observation records a rational allowance $\epsilon$ and four
rational coordinates $y_i$. Compatibility is the stated coordinatewise error
contract, together with the admitted unit bounds and allowance conditions.
All observations in one history must fit \emph{the same} latent state.

For a well-formed finite history, compatible coordinates lie in a rectangular
box with endpoints
\[
 L_i=\max\bigl(0,\max_j(y_{ij}-\epsilon_j)\bigr),\qquad
 U_i=\min\bigl(1,\min_j(y_{ij}+\epsilon_j)\bigr).
\]
The empty history starts with $[0,1]^4$. Malformed allowances and empty
intersections are handled by the implemented feasibility checks; the formula
above is not permission to ignore input validity.
\leanref{Nullivance.Recognition.summarizeProbes_exact}
\leanref{Nullivance.Recognition.historyFeasible_iff_exists}

\begin{theorem}[Exact finite classification]
Every finite input history receives precisely one of: invalid (no compatible
state), affirmed (nonempty and all compatible states quasivant), refuted
(nonempty and all compatible states nonquasivant), or undetermined (compatible
states of both kinds).
\end{theorem}
\begin{proof}
For each coordinate retain the two interval endpoints and its distinguished
marker: zero for intensity, $1/2$ for structure. This gives $3^4=81$ candidate
entries, not necessarily distinct or feasible. Search only candidates passing
all bounds and the requested truth value of $\Qv$. Every real compatible state
has a candidate preserving the four equalities that determine $\Qv$; conversely
each accepted candidate embeds into an admitted state. The two existence
flags determine the four disjoint cases. Feasibility prevents vacuous
affirmation on inconsistent data.
\end{proof}
\leanref{Nullivance.Recognition.boxWitnessCandidates_represent}
\leanref{Nullivance.Recognition.exportHistoryWitness_realizes}
\leanref{Nullivance.Recognition.classifyHistory_complete}

\begin{example}[Reading the four verdicts]
Consider boxes with both structure coordinates in $[0,1/4]$, so neither can
be neutral. If both intensity intervals are $[0,0]$, the box is nonempty and
every member satisfies $\Qv$: the verdict is affirmed. If the positive
intensity interval is $[1/10,1/5]$, every member violates $\Qv$: refuted.
If both intensities instead range over $[0,1/10]$, zero intensities and positive
intensities are both possible: undetermined. Finally, contradictory constraints
requiring one coordinate to be at most $1/4$ and at least $3/4$ have no common
state: invalid. These elementary box examples illustrate the classifier's
specification, not empirical measurements.
\end{example}
An undetermined verdict therefore contains positive mathematical information:
both kinds of state remain possible. An invalid verdict identifies failure of
the common-state/error contract and must not be interpreted as evidence for
every property. These four verdicts describe sets of compatible latent states;
they are not a renaming of the logical values N, T, F, B.

The rational witness is actual exported data, not a claim that an arbitrary
real hidden state can be computed. Completeness here concerns this predicate
and independent coordinate boxes. Correlated nonlinear constraints would
require another representation theorem. Permuting or duplicating observations
does not alter the compatible set; the artifact proves the corresponding
implementation invariances.

\begin{example}[Combining ambiguous records]
With allowance $1/10$, readouts $((0,0),(2/5,0))$ and
$((0,0),(3/5,0))$ are separately undetermined. Together their positive-structure
interval is exactly $[1/2,1/2]$. Every common state has a neutral positive
structure, so the combined verdict is refuted, not invalid.
\end{example}
\leanref{Nullivance.Recognition.joint_neutrality_refutation_regression}
The example illustrates intersection of evidence about one state; it does not
model an unannounced change of state between readings.

\section{Infinite histories and small certificates}
For the same scalar, closed, independent-coordinate model, a single state
fits a countable stream iff every finite prefix is feasible. In fact it is
enough that every pair of records has a common fitting state.
\begin{proof}[Argument for the compatibility result]
Pairwise feasibility implies each lower endpoint is at most every upper
endpoint in the same coordinate. Take the supremum of all lower endpoints.
The set is nonempty and bounded above by an upper endpoint from record zero.
Its supremum lies between all the lower and upper endpoints and in $[0,1]$.
The four coordinates assemble one admitted state fitting the entire stream.
The reverse directions follow by restricting a fitting state.
\end{proof}
\leanref{Nullivance.Recognition.stream_exists_iff_all_prefixes_feasible}
\leanref{Nullivance.Recognition.stream_exists_iff_pairwise_feasible}
Thus inconsistency has a certificate of at most two observations. Two is
sharp: exact positive-intensity readings zero and one are individually feasible
but jointly inconsistent. This is a rectangular-constraint result, not a new
unrestricted compactness theorem. Certificate size does not bound where the
records occur or how long the observer waits.
\leanref{Nullivance.Recognition.stream_infeasible_iff_two_observation_certificate}
\leanref{Nullivance.Recognition.two_observation_bound_is_sharp}

Compatibility compactness must not be confused with finite certification of
every universal property. An infinite shrinking history can force a boundary
value while every finite prefix still admits alternatives. The development
therefore preserves explicit non-stopping examples, instead of assuming that
infinite evidence always becomes finite evidence.

The certificate branch makes these distinctions operational. Endpoint summaries
carry the original observations responsible for extremal bounds. Conflict,
affirmative, and refutation certificates are checked against source positions.
Feasibility is retained for non-vacuous positive or negative conclusions.
An actual source index is more informative than an unattributed value: a
numerically correct endpoint alone does not authenticate its claimed origin.
The certificate and memory refinements develop extraction, indexed provenance, joint
pipelines, cursor-field costs, and short-circuit equality refinement. Each
cost theorem refers to its own instrumented implementation; output equivalence
alone is not used to transfer an unproved runtime bound.

\section{Imported memory, trusted continuation, and serialization}
\label{sec:memory}
A saved cursor includes its next position and a sourced summary. On import,
the externally supplied object is not assumed to match the underlying records.
The checked entry recomputes the prefix summary and compares it with the
imported value. It returns a proof-carrying handle precisely when
\nolinkurl{CursorMatches} holds. This invariant concerns equality to the summary
of the designated source prefix; it is not cryptographic proof that an external
provider reported physical truth.
\leanref{Nullivance.Recognition.importTrusted_isSome}
\leanref{Nullivance.Recognition.importTrusted_output}

\begin{proposition}[Counted import]
For cursor position $n$, the specified counted importer equals the checked
import result together with $8n$ rational-order nodes and one whole-box equality
call, whether it accepts or rejects.
\end{proposition}
\begin{proof}
The counted source-prefix check has the same Boolean result as the ordinary
check and the displayed counters. The importer branches on that Boolean:
acceptance packages the same cursor with the established matching proof;
rejection returns no handle. Neither branch changes the counters. Substitution
of the prefix-check specification therefore gives the result in both cases.
The proof field is a logical guarantee, not additional serialized evidence.
\end{proof}
\leanref{Nullivance.Recognition.importTrustedCounted_spec}
The unit ``one equality call'' is deliberately not one machine instruction:
later refinements count fields and short-circuit behavior separately. Once
accepted, an in-process proof-carrying continuation preserves the invariant
without rescanning the old prefix at every step. The stream is fixed during
these proofs. Mutating the source after validation is outside the contract.

A finite digest cannot universally replace exact checking: any map from the
relevant infinite family into a finite digest space collides. The formal
counterexample includes feasible histories, so the obstruction is not merely
malformed data. This is an impossibility of \emph{unconditional exact inference}
from a finite digest, not a refutation of probabilistic cryptographic security
under a specified threat model.
\leanref{Nullivance.Recognition.finite_digest_cannot_justify_import}

The serialization layer specifies session lifecycle and checked byte/envelope
boundaries. A Lean structure carrying a proof is not by itself a sealed runtime
capability: proof fields are erased. Reimporting raw data requires checking
again or an explicitly specified trusted-storage assumption. The later cache,
checkpoint, and typed-phase imports likewise validate the relevant semantic
invariants rather than treating serialization success as semantic correctness.
The project does not yet prove a general secure distributed-storage system.
\leanref{Nullivance.Recognition.decodeCursorEnvelope_encode}
\leanref{Nullivance.Recognition.restoreActiveEnvelopeCounted_success}
The separate threshold-endpoint cache validator also has an exact acceptance
specification relative to rebuilding from supplied threshold data. Its accepted
path costs $2n+11$ selected operations when $n<2$, and $2n+18$ otherwise.
The model charges list-constructor inspections, stage-tag inspections, and
Boolean gates; its comparison circuit is charged without short-circuit savings.
This is a different cost unit from
the rational-order nodes of sourced-history import. Acceptance does not imply
that arbitrary supplied threshold data are realizable by an underlying state,
nor that two full histories are equal.
\leanref{Nullivance.Recognition.checkThresholdCache_spec}
\leanref{Nullivance.Recognition.checkThresholdCacheCounted_accepted}
\leanref{Nullivance.Recognition.checked_cache_nonrealizable_counterexample}

\section{Threshold observations and finite information limits}
There are two distinct information models: exact or bounded-error scalar
readings, and observations restricted to threshold bits. Finitely many positive
thresholds cannot in general distinguish exact zero from every sufficiently
small positive intensity. An appropriate infinite countable threshold family
can determine zero as a mathematical property of the entire signature, but
this does not produce a finite algorithm affirming zero in every case.
The quasivance predicate additionally requires both structures to be nonneutral.
An intensity-only result cannot settle that conjunction.
\leanref{Nullivance.Recognition.finite_positive_thresholds_not_recognizable}
\leanref{Nullivance.Recognition.countableThreshold_recognizes_quasivance}

The threshold analysis distinguishes init-only and partial-probe
losses, finite versus countable threshold signatures, margin assumptions,
neutrality boundaries, finite refutation, and endpoint compression.
The finite-prefix classifier retains its distinction between invalid,
affirmed, refuted, and undetermined. A below-threshold status is not equality
to zero and is not automatically affirmation of quasivance. Conversely an
above-positive-threshold certificate refutes quasivance by witnessing positive
intensity in at least one channel.

The streaming implementation connects these information limits to supplied-list costs,
incremental summaries, checked imported caches, first-hit stopping,
consumed-prefix accounting, chunk composition, checkpoint continuation,
session ledgers, ordered delivery, and explicit precision deadlines.
``Eventually delivered'' is a hypothesis about records received by the
algorithm, not a theorem that a physical sensor will supply them.
Selective refresh preserves resolved results only under the proved refinement
and feasibility conditions. Successful chunk composition does not authorize
reordering records or skipping the source-prefix checks.

\section{Adaptive resolution and the exact boundary}
\label{sec:boundary}
Let $I(s)=\max(\alpha_+,\alpha_-)$. A finite list of positive rational
thresholds is separated from this value when $I(s)\ne\tau$ for each listed
$\tau$. Such a finite separated list admits a common positive rational margin.
Existence of that margin is a proof fact, not a decoder extracting it from
arbitrary observations. With the specified delivery of arbitrarily fine
precision, the bounded reference search eventually finds a prefix resolving
all requested thresholds. It does not receive the latent state, margin, or
delivery schedule as executable inputs.
\leanref{Nullivance.Recognition.separated_stream_adaptive_resolution}

The search starting at prefix length $a$ with fuel $b$ tests lengths
$a,\ldots,a+b-1$. It returns the first ready prefix or $\none$. Here
$\none$ means no success within that window, not a negative theorem about the
latent state. Zero fuel tests nothing. Empty threshold lists still require
feasibility. The executable refinements implement this reference search as
incremental, whole-run and read-on-demand replay, including shortage, resume,
typed phase import, and compact initial-prefix summaries. These refinements
preserve the reference first-success result and separately account for the
operations in their instrumented implementations.
The compact initialization summary retains eight rational endpoints and at
most eight occupied observation/index slots. This bound is attained for that
representation; it is not a lower bound for every algorithm. Equal sourced
summaries give equal future phases for fixed requests, fuel, and a common
ordered suffix. They do not determine the original history: different records
can yield the same summary. Thus a correct resumable memory is not a lossless
archive of research data.
\leanref{Nullivance.Recognition.compactReplay_extension}
\leanref{Nullivance.Recognition.bootSources_bound}
\leanref{Nullivance.Recognition.bootSources_eight_attained}
\leanref{Nullivance.Recognition.no_bootSummary_history_inverse}

\subsection{A necessary and sufficient finite boundary criterion}
For one observation $r_j=(\epsilon_j,y_j)$ define its intensity floor
\[
 f(r_j)=\max(y_{+,j}-\epsilon_j,\ y_{-,j}-\epsilon_j).
\]
Every fitting state's maximum intensity is at least this floor. Suppose a
common fitting state has $I(s)=\tau>0$. Then a finite history resolves this
threshold iff some observed floor equals $\tau$.
\leanref{Nullivance.Recognition.boundary_ready_iff_contact}
\begin{proof}
The feasible box certifies ``at least $\tau$'' precisely when one lower intensity
bound reaches $\tau$. Its initial value zero cannot do so because $\tau>0$.
Every later lower bound is the maximum of prior lower bounds and observed
lower endpoints. Each such endpoint is at most the fitting intensity, hence
the largest floor is at most $I(s)=\tau$. Reaching the threshold therefore
means equality at an observed floor. A below verdict is excluded by the
given fitting boundary state. If there is no contact, the simultaneous lower
corner of the feasible box has both intensities below $\tau$ and is itself
a fitting state. The finite history then admits both a boundary state and
a below-threshold alternative.
\end{proof}
\leanref{Nullivance.Recognition.boundary_contact_or_alternative}

For a threshold in $(0,1]$, constant centered intensity readings with allowances
$1/(j+1)$ can force $I=\tau$ in the whole infinite stream while every finite
search times out: each floor remains strictly below $\tau$. Arbitrarily fine
precision alone is therefore insufficient for universal finite boundary
resolution. Conversely the reading $y_+=3/4$ with allowance $1/4$ has floor
$1/2$ and certifies at least $1/2$. Exact zero noise is not necessary.
\leanref{Nullivance.Recognition.centeredBoundary_timeout}
\leanref{Nullivance.Recognition.noisy_contact_regression}
The successful noisy reading certifies an inequality, not exact equality to
$1/2$. Equality is an assumption of the boundary analysis, not a conclusion
silently inferred from that reading.

\section{Finite generators, equivalence, and counted cycle bounds}
A contact generator consists of a start state, a deterministic next-state
function, and a Boolean mark. With a finite state space of size $D$, if a
mark is ever reached then one is reached before step $D$. A bounded scanner
returns the first marked position; exhaustion is a permanent negative only
when the applicable cutoff theorem has been established.

Generator descriptions can provide information absent from sampled traces.
The development includes finite-product trace comparison, checked finite
relations, symbolic simulation relations, and residue-counter encodings.
A symbolic relation certificate can establish equality of all output traces,
but existence of such a relation is not an algorithm discovering it.
Similarly an affine-residue checker is complete for the proposed affine
invariant it checks, not for every pair of trace-equivalent generators.
Unreachable states can invalidate that invariant without changing reachable
output behavior.

For a residue counter with modulus $P>0$ and stride $s$, the reachable cycle
length is $c=P/\gcd(P,s)$. In particular stride zero gives $c=1$, not zero.
Two such traces with cycle lengths $c,d$ can be compared through
$L=\operatorname{lcm}(c,d)$, bounded by the earlier product cutoff.
The counted comparison preserves the exact earliest mismatch.
The cost analysis explicitly instruments a Euclidean algorithm for
setup, proves a logarithmic remainder-step bound by two-step halving, and
avoids setup entirely when the first marks already differ.

For the specified setup, if
$S=\log_2 P+\log_2 Q+\log_2 c$ with natural-number logarithms,
the proved selected-work envelope is $2S+13+2L$.
This counts the named remainder/division/product/mark/transition operations.
It is not the bit complexity of Euclid or the runtime of any uninstrumented
built-in gcd implementation. A large-number division has not been proved to
cost one processor step. Equivalence to a faster-looking algorithm does not
change this scope.
\leanref{Nullivance.Recognition.compareResidueCycles_exact}
\leanref{Nullivance.Recognition.euclidCounted_log_bound}
\leanref{Nullivance.Recognition.preparedCycleWork_log_bound}
\leanref{Nullivance.Recognition.guardedCycles_result}

\section{Periodicity without a bound is still insufficient}
\label{sec:negative}
Here periodicity refers to the Boolean contact pattern,
not necessarily the full rational observation stream. A residue-generator
description enables a cutoff calculation. Merely promising some unknown
period to an observation-only algorithm does not supply that calculation.

\begin{theorem}[Finite negative query obstruction]
There is no deterministic observation-only strategy that finitely returns a
sound negative answer exactly on every never-contacting stream in the
specified periodic encoding family, without being supplied a period bound.
Queries may be adaptive, repeated, or out of order.
\end{theorem}
\begin{proof}[Finite-transcript argument]
Run the strategy on the silent encoded counter. Any terminating negative run
has queried finitely many indices. Choose $N$ larger than all of them. A
late-contact periodic counter with cycle $N+1$ agrees on every queried
observation but first contacts at $N$. Determinism preserves the same
transcript and negative answer, contradicting soundness on this later-contact
stream. The encoding maintains a common admitted scalar state and the stated
precision properties, so the counterexample is inside the modeled class.
\end{proof}
\leanref{Nullivance.Recognition.periodicQuery_admissible_collision}
\leanref{Nullivance.Recognition.no_finite_negative_periodic_recognizer}
This is not a theorem about randomized algorithms, arbitrary source-code
analysis, or uncomputability of every particular generator. It also does not
say positive contacts are undiscoverable: a finite observed contact is a
positive certificate. The obstruction concerns universally conclusive finite
\emph{negative} observation-only decisions under the stated interface.

\section{Decision with a certified period bound}
\label{sec:149}
Define $C(n)$ by $f(r_n)=\tau$. The new assumption is
\[
 \mathsf{ContactPeriodBound}(r,\tau,B):\quad
 \exists P>0,\quad P\le B\ \land\ \forall n,
 \bigl(C(n+P)\leftrightarrow C(n)\bigr).
\]
This is an all-index semantic promise, not a Boolean test performed on the
first $B$ samples. The executable scanner receives $r,\tau,B$, not the
period or generator description. It checks indices $0,\ldots,B-1$ in order
and stops at the first contact.

\begin{theorem}[Certified cutoff and scanner specification]
Under the period-bound promise, contact ever occurs iff it occurs at some
$k<B$. The scanner returns $\some(k)$ exactly when $k<B$, $C(k)$, and
no earlier index contacts; it returns $\none$ exactly when contact never
occurs. The first-hit specification needs no periodicity assumption; the
permanent-negative specification does.
\end{theorem}
\begin{proof}
If a contact exists, take its least index $m$. If $m\ge P$, periodicity yields
a contact at $m-P<m$, a contradiction. Hence $m<P\le B$. The reverse
implication is immediate. The bounded scanner's first-hit and finite-exhaustion
specifications follow from its recursive mark test. Composing finite
exhaustion with the cutoff yields the permanent-negative statement.
\end{proof}
\leanref{Nullivance.Recognition.contactPeriodBound_cutoff}
\leanref{Nullivance.Recognition.scanObservedContact_some}
\leanref{Nullivance.Recognition.scanObservedContact_permanent}

If all observations fit one admitted scalar state with $I(s)=\tau>0$, the
negative answer is equivalent to the canonical threshold procedure never
finding a ready prefix. A first contact at $k$ gives its first successful
prefix $k+1$, tested with budget $B+1$ because the empty prefix is included.
No precision-delivery assumption is required for this exact-boundary bridge.
The negative result means ``no finite boundary certificate in this procedure,''
not ``the intensity is below $\tau$'': the theorem assumes equality.
\leanref{Nullivance.Recognition.scanObservedContact_boundary}
\leanref{Nullivance.Recognition.scanObservedContact_first_prefix}

\begin{theorem}[Exact selected-operation counts]
The counted scanner returns the same result as the scanner. At a first hit
$k$, it charges $(k+1,k)$ contact tests and successor transitions. At exhaustion
it charges $(B,B)$. Both counters are at most $B$.
\end{theorem}
\begin{proof}
Instantiate the counted generator scanner with the actual rational floor
equality as its mark operation and natural successor as its step. Each failed
test advances once, including the final failing transition on exhaustion;
a successful test does not advance. The exact recurrence gives the two
cases and, since $k<B$, the bounds.
\end{proof}
\leanref{Nullivance.Recognition.scanObservedContactCounted_exact}
\leanref{Nullivance.Recognition.scanObservedContactCounted_boundary}
This charges the operations of the actual wrapper, not a cost inferred merely
from output equivalence. It excludes sample production, rational bit arithmetic,
allocation, and compiler-level callback behavior. Computability is relative to
an evaluable stream; an arbitrary mathematical function is not a hardware
sensor driver.

The late-contact family proves sharpness of this leading-prefix bound:
with period $N+1$, $N$ observations return no hit, whereas $N+1$ observations
first hit at $N$. This is not an adaptive-query lower bound for the bounded
class. Budget zero is executable and returns no hit with zero work, but it
cannot satisfy the positive-period promise. Testing a few samples for apparent
periodicity does not justify promoting exhaustion to a permanent negative.
\leanref{Nullivance.Recognition.scanObservedContact_cutoff_sharp}
\leanref{Nullivance.Recognition.contactPeriodBound_positive}

\section{Formal verification and reproducibility}
\label{sec:repro}
The verification uses Lean 4 \cite{demoura2021lean}. The companion artifact
pins the mathematical baseline, Lean toolchain,
mathlib commit, and all dependency revisions in the Lake manifest. Its
\nolinkurl{SOURCE_MANIFEST.sha256} covers packaged files; a separate digest
covers the source ZIP. Hashes detect byte changes relative to that manifest;
they are not proofs of scientific truth or external signing credentials.

The verification script runs the library build, both research validation
entry points, and the overview and companion declaration/type/dependency audits.
It rejects proof holes and nonstandard axioms in the report's theorem anchors.
The complete validation logs remain available for inspection. Legacy executable
regressions that use a native decision bridge must not be confused with the
structural mathematical proofs: any such dependency is disclosed in the
validation inventory rather than erased from the record. The latest
recognition boundary results use kernel-checked proofs and standard Lean
logical dependencies only.

The expected standard logical dependencies are \nolinkurl{propext},
\nolinkurl{Classical.choice}, and \nolinkurl{Quot.sound}, or subsets.
``No added axioms'' does not mean that classical mathematics has no foundational
assumptions. Axiom audits inspect the dependency closure of a declaration;
source-token scans alone would not suffice.

The theorem-anchor audit prints full elaborated types. This permits checking
that, for example, a theorem is conditional on \nolinkurl{StreamFits} and
\nolinkurl{ContactPeriodBound}, rather than simply checking that its name
exists. The publication review additionally inspects definitions, algorithm
inputs, cost units, and counterexamples. The associated audit report states
which checks were actually executed and does not label a cache-assisted
extracted build as a cold dependency rebuild.

To reproduce, install the pinned toolchain and LaTeX prerequisites, use the
packaged instructions, and run from \nolinkurl{Nullivance/}:
\begin{verbatim}
lake build
lake env lean ResearchValidation.lean
lake env lean RecognitionValidation.lean
lake env lean PublicationValidation.lean
lake env lean CompanionValidation.lean
\end{verbatim}
The last two files are generated from the theorem anchors in all three reports.
The package also contains the manuscript source, bibliography, and generated
catalog, so rebuilding the report does not require private chat records or
local drive paths. Source-only reproduction may need network access to obtain
the pinned dependencies. No credentials, agent conversation logs, or worktree
administration files belong in this publication artifact.

\section{Discussion, limitations, and conclusion}
\subsection{What the results mean}
The logical and recognition layers answer different questions. The first
controls consequence from support values. The second controls conclusions
about a hidden state given an observation contract. Their connection requires
explicit assumptions; no theorem turns arbitrary reports into factive knowledge.
In particular, the initial-observation counterexample identifies a failure of
an interface, whereas the scalar probe construction specifies how a stronger
interface can remove that failure.

The stopping results isolate another distinction. Infinitely improving precision
can determine a boundary value without any finite history certifying it. A
period bound changes the available specification, not merely the speed of
the same uninformed search. The positive and negative theorems are consistent
because their input contracts differ. This separation is useful when designing
software that must report uncertainty instead of silently treating timeout as
a negative answer.

\subsection{Outstanding obligations}
A central outstanding obligation is to obtain or validate a
usable period or tail certificate from an independently justified application
source. It cannot in general be manufactured from the same finite observations
whose negative result needs justification. Generator-description verification
and assumptions about a physical generating process are different tasks.

Other boundaries remain: unrestricted infinite-domain quantified completeness;
validity across all assignments beyond the supplied-assignment equality
interfaces; arbitrary structural dimensions and changing-state observations;
concrete adequate evidence maps; correlated observation errors; combined
serialization and application-security contracts; bit complexity and physical
resource bounds; and empirical validation of any cognitive interpretation.
Some narrower obligations have already been closed by later records. This
list therefore describes current scope rather than repeating every historic
``next step'' as if it were still open.

Potential applications include checking whether an observation interface
supports a requested distinction, returning explicit alternatives when it
does not, validating cached summaries against source data, and reporting
certified stopping only under justified assumptions. These are mathematically
motivated opportunities, not measured superiority over existing systems.
No benchmark or human/AI cognitive experiment is reported here.

\subsection{Conclusion}
The main research contribution of this collection is a traceable formal chain:
precise semantics; scoped completeness and equality algorithms; observation
loss counterexamples; exact finite recognition under a specified probe model;
source-backed, checked implementations; and a matched impossibility/positive
cutoff pair for periodic boundary contact. The introductory comparison identifies
relevant established ideas but does not establish priority for the full
construction. We do not claim that these general
ideas originated with this artifact or that the complete problem of
recognition has been solved.

\section*{Availability and author declarations}
The project repository is
\url{https://github.com/lamtung0487-droid/nullivance-logic}.
The attached source archive identifies this snapshot independently of whether
that repository has subsequently advanced. The artifact is distributed under
the project's CC BY 4.0 license. The author contact is given on the title page.
The version family is \url{https://doi.org/10.5281/zenodo.18309848};
the record labeled \PublicationVersion\ identifies this collection specifically.
An archival deposit is not external peer review or empirical certification.
No affiliation, funding, or competing-interest statement is inferred from
the software artifact.

\appendix
\section{Collection identity and relation to earlier editions}
This edition changes exposition, bibliography, version labeling, and verification
tooling, not the mathematical baseline through DR-0149:
\nolinkurl{48bd5f4260fd7982176d7a2060be4e956e1fd49c}.
Lean and mathlib remain pinned to 4.32.1. Exact packaged bytes are identified by
the source manifest and successful verification receipt, not by the baseline
commit alone.

\begin{center}\small
\begin{tabular}{p{.19\linewidth}p{.70\linewidth}}
\toprule Edition & How it should be read\\\midrule
Original version 1 & Historical motivation and exploratory formulation;
not evidence that every philosophical or physical interpretation was proved.\\
v0.7.1 & Archived propositional and finite-domain papers with their original
formalization status; preserved at their original record.\\
DR0149-r2 & Prior synthesis and reproducible artifact. Its bundled core PDF
included later formalization updates despite retaining the old cover label.\\
DR0149-r3 & This consistently labeled three-part collection. Part II retains
those updates and identifies them; Part III retains its mathematics with revised
presentation. The overview is rewritten around questions, proofs, and limits.\\
\bottomrule
\end{tabular}
\end{center}
Current formal definitions take precedence for current claims. In particular,
exact zero is not a small positive intensity; independent support coordinates
are not complementary probabilities; a proposed latent structure is not an
observed physical cause. Earlier broad comparisons with classical logic do not
establish limitations of classical mathematics. These clarifications preserve
historical provenance without elevating exploratory claims to theorems.

\section{Complete chronological coverage through DR-0149}
\label{app:catalog}
The following catalog preserves every design-record identifier and title in
the source baseline. It is a coverage index, not 149 claims of novelty or a
claim that every historical proposed statement was proved. Records include
definitions, counterexamples, repaired proof routes, algorithms, and cost
refinements. The accompanying machine-readable catalog lists original record
paths, SHA-256 identities, and explicitly mentioned Lean files when present.
Historical verification counts in those records are not substituted for the
fresh publication checks.
\small
\begin{longtable}{p{.14\linewidth}p{.79\linewidth}}
\toprule Record & Topic\\\midrule\endhead
DR-0001 & Propositional language of NPL\\
DR-0002 & Two-channel continuous semantics and connective clauses\\
DR-0003 & Threshold projection, meta-signs, and the consequence relation\\
DR-0004 & Lean encoding choices and the continuous-layer gap\\
DR-0005 & The four-signed tableau calculus is the canonical proof system\\
DR-0006 & Formalizing the generative tier (Tier 1)\\
DR-0007 & Finite-domain quantified NPL first\\
DR-0008 & Finite-domain quantified tableau first\\
DR-0009 & Equality closure repair for finite-domain quantified tableau\\
DR-0010 & Finite grounding bridge for quantified completeness\\
DR-0011 & Full recursive extensional finite-domain tableau\\
DR-0012 & Constrained grounding completeness for finite-domain tableau\\
DR-0013 & Core extensional tableau and replay trace\\
DR-0014 & Fixed-signature layer for finite quantified NPL\\
DR-0015 & Epistemic statuses for definitions and conventions\\
DR-0016 & Exact continuous-semantics encoding bridges\\
DR-0017 & Verified operational reference search\\
DR-0018 & Progress-scheduler independence without fairness\\
DR-0019 & Refutation of the head-sensitive replay bridge\\
DR-0020 & Membership-selecting replay certificate\\
DR-0021 & Arbitrary-domain quantified semantic foundation\\
DR-0022 & Collision-sensitive cost and finite-domain certificates\\
DR-0023 & Open equality: assignment partitions and fresh capacity\\
DR-0024 & Executable free-assignment encoding\\
DR-0025 & Compact open orbit initialization\\
DR-0026 & Exact finite open equality and checked admission\\
DR-0027 & Recognition, observation loss, and extension criteria\\
DR-0028 & Proposed scalar active probes and their boundaries\\
DR-0029 & Bounded-noise recognition and a rational decoder\\
DR-0030 & Sound abstaining probe classifier without a margin promise\\
DR-0031 & Exact feasibility and completeness of the abstaining classifier\\
DR-0032 & Executable rational recognition witnesses\\
DR-0033 & Exact same-state observation histories\\
DR-0034 & Complete affirmation for finite observation histories\\
DR-0035 & Complete history classification and executable witnesses\\
DR-0036 & Infinite histories need not yield finite determination\\
DR-0037 & Compactness and sharp two-observation inconsistency certificates\\
DR-0038 & Executable small conflict certificates\\
DR-0039 & Sourced linear history summaries and indexed conflict certificates\\
DR-0040 & Finite recognition: attained boundaries versus limiting information\\
DR-0041 & Persistence of finite recognition certificates\\
DR-0042 & Executable bounded recognition-certificate search\\
DR-0043 & Earliest certificate and exact abstract test count\\
DR-0044 & Incremental, resumable recognition-certificate search\\
DR-0045 & Exact finite refutation and two-record certificates\\
DR-0046 & Endpoint classifier and composed comparison budget\\
DR-0047 & Executable indexed refutation certificates\\
DR-0048 & Refutation export from saved search state\\
DR-0049 & Search plus certificate-export comparison budget\\
DR-0050 & Authenticated index-only source deduplication\\
DR-0051 & Finite stopping from a positive-intensity margin\\
DR-0052 & Computable reciprocal-error stopping budget\\
DR-0053 & Affirmative stopping with attained clamps and theta separation\\
DR-0054 & Authenticated finite affirmative certificates\\
DR-0055 & Affirmative export from a saved search cursor\\
DR-0056 & Counted affirmative export and search composition\\
DR-0057 & Reusing affirmative guard comparisons\\
DR-0058 & Authenticated index-only affirmative export\\
DR-0059 & Computable authentication of supplied source boxes\\
DR-0060 & Checked continuation from an imported cursor\\
DR-0061 & Affirmative certificates after checked continuation\\
DR-0062 & Refutation certificates after checked continuation\\
DR-0063 & Exact-one certificates and joint operational outcomes\\
DR-0064 & Validation-inclusive continuation counters\\
DR-0065 & Counted index-only refutation and joint pipeline\\
DR-0066 & Sharing final-export consistency checks\\
DR-0067 & Field-level charged-node budget for cursor validation\\
DR-0068 & Thread field-level cursor charges through the joint pipeline\\
DR-0069 & Explicit short-circuit equality for source records\\
DR-0070 & Source short-circuit comparison through checked joint certificates\\
DR-0071 & No source short-circuit saving on an accepted cursor\\
DR-0072 & Finite digests alone cannot authenticate every saved cursor\\
DR-0073 & Finite-digest failure persists for feasible alpha/Theta histories\\
DR-0074 & Proof-carrying in-process cursor continuation\\
DR-0075 & Checked entry from untrusted saved cursors\\
DR-0076 & Terminal-safe trusted sessions and cumulative schedule bound\\
DR-0077 & Checked import followed by a trusted finite schedule\\
DR-0078 & Active-only checkpoint and revalidation boundary\\
DR-0079 & Canonical typed-byte checkpoint codec\\
DR-0080 & Versioned ByteArray checkpoint envelope and size gate\\
DR-0081 & A first-stage numerical bound for checkpoint decoding\\
DR-0082 & A non-vacuous bridge from finite histories to factive knowledge\\
DR-0083 & Conditional evidence-to-knowledge bridge and two obstructions\\
DR-0084 & Init-only FOUR reports cannot exactly detect quasivance\\
DR-0085 & Each of the two existing active probes is individually insufficient\\
DR-0086 & Finitely many positive FOUR thresholds remain insufficient\\
DR-0087 & An infinite shrinking-threshold trace identifies zero intensity\\
DR-0088 & Whole countable threshold signature recognizes scalar quasivance\\
DR-0089 & Three threshold levels suffice under an explicit scalar margin\\
DR-0090 & Finite one-sided intensity refutation without a uniform margin\\
DR-0091 & Finite witnesses of scalar structure nonneutrality\\
DR-0092 & No finite prefix affirms scalar quasivance\\
DR-0093 & A neutral nonquasivant boundary cannot be refuted by finite prefixes\\
DR-0094 & Every zero-intensity scalar prefix admits a quasivant state\\
DR-0095 & Exact eventual finite-prefix refutation criterion\\
DR-0096 & The zero-length full prefix already has a refutation bit\\
DR-0097 & Exact refutation criterion for every full finite prefix\\
DR-0098 & Compressing full-prefix intensity refutation\\
DR-0099 & Single-threshold summary of finite-prefix refutation\\
DR-0100 & Supplied-list refutation and exact traversal cost\\
DR-0101 & Incremental threshold endpoint and update cost\\
DR-0102 & Imported threshold-cache validation and replay cost\\
DR-0103 & Proof-carrying threshold import and preserved update invariant\\
DR-0104 & Checked import followed by a finite online trace\\
DR-0105 & First-hit threshold stopping and scalar persistence\\
DR-0106 & Checked stopping, chunk continuation and explicit offset cost\\
DR-0107 & Retained consumption and repeated-chunk cost\\
DR-0108 & Checked import followed by certified chunk stopping\\
DR-0109 & Validation of stopping-checkpoint metadata\\
DR-0110 & Continuing an accepted stopping certificate\\
DR-0111 & Finite continuation sessions and a separate cost ledger\\
DR-0112 & Ordered delivery bridges finite sessions to eventual refutation\\
DR-0113 & Explicit rational-margin observation and delivery budgets\\
DR-0114 & Feasibility-guarded threshold certificates from error histories\\
DR-0115 & Searching partial intensity certificates without false completion\\
DR-0116 & Lazy certificates and safe replay after observation refinement\\
DR-0117 & Shared materialized summaries for partial threshold replay\\
DR-0118 & Length-carrying summaries and finite-session cost\\
DR-0119 & Reusing resolved thresholds under feasible refinement\\
DR-0120 & Counted selective-cache initialization and finite sessions\\
DR-0121 & Checked import of the combined selective cache\\
DR-0122 & Finite unknown-exit potential and qualitative stabilization\\
DR-0123 & Precision-delivery deadlines for separated intensity thresholds\\
DR-0124 & Adaptive prefix resolution without a prior margin\\
DR-0125 & Incremental all-threshold resolution search\\
DR-0126 & Exact cache and cost composition across search budgets\\
DR-0127 & Checking an imported resolution run against finite source data\\
DR-0128 & Replaying only the observations a completed search needs\\
DR-0129 & Charged work of a replay that runs out of observations\\
DR-0130 & Resume replay after an observation shortage\\
DR-0131 & What must be checked before resuming an imported pending phase?\\
DR-0132 & Receive the starting prefix incrementally\\
DR-0133 & A sufficient sourced summary for incomplete initialization\\
DR-0134 & When does a finite history decide an exact intensity boundary?\\
DR-0135 & Exact eventual stopping for a finite list of intensity thresholds\\
DR-0136 & No finite observational certificate of permanent non-resolution in general\\
DR-0137 & A sharp cutoff given a finite generator of boundary contacts\\
DR-0138 & Finite validation of a candidate contact generator\\
DR-0139 & Checkable finite relation certificates for generator traces\\
DR-0140 & Proof-carrying symbolic relations for infinite counter generators\\
DR-0141 & An executable checker for affine relations of periodic counters\\
DR-0142 & Exact trace decision through a finite residue quotient\\
DR-0143 & Stride-sensitive synchronized cycle cutoff\\
DR-0144 & Verified arithmetic setup and scanner work\\
DR-0145 & Logarithmic Euclidean remainder-count bounds\\
DR-0146 & Skip cycle preparation on an immediate disagreement\\
DR-0147 & Periodic contact decision and the unknown-period prefix limit\\
DR-0148 & Adaptive querying does not remove the unknown-period limit\\
DR-0149 & Observation-only decision under a certified period bound\\
\bottomrule
\end{longtable}
\normalsize


\bibliographystyle{plain}
\bibliography{snapshot-references}
\end{document}
